\documentclass{article}
\usepackage{amsmath,amssymb,amsthm}
\usepackage{algorithm}
\usepackage[noend]{algpseudocode}
\usepackage{graphicx}
\usepackage{hyperref}
\usepackage{enumitem}
\newtheorem{theorem}{Theorem}
\newtheorem{lemma}{Lemma}
\newtheorem{assumption}{Assumption}
\newtheorem{remark}{Remark}
\newcommand{\E}{\mathbb{E}}
\newcommand{\R}{\mathbb{R}}
\newcommand{\norm}[1]{\left\|#1\right\|}
\newcommand{\inner}[1]{\left\langle#1\right\rangle}
\newcommand{\topk}[1]{\mathcal{S}_K\!\left(#1\right)}
\begin{document}

\section*{Top-$K$ Error‑Feedback SGD (EF‑TopK‑SGD)}

\section*{Mathematical Formulation}
Let $f:\R^d\rightarrow\R$ be a (possibly non‑convex) smooth objective and let $\{\xi_t\}_{t\ge 0}$ be i.i.d.\ data samples.  
Define the stochastic gradient
\[
g_t \;:=\; \nabla f(w_t;\xi_t)\in\R^d .
\]

\paragraph{Error‑feedback variables.}
\begin{itemize}
    \item $e_t\in\R^d$ stores the accumulated compression error up to iteration $t$.
    \item The \emph{corrected} gradient is $u_t = g_t+e_t$.
\end{itemize}

\paragraph{Top‑$K$ compressor.}
For $K\in\{1,\dots,d\}$ define the operator $\topk{\cdot}:\R^d\to\R^d$ as
\[
\bigl[\topk{x}\bigr]_i=
\begin{cases}
x_i, & \text{if } i\in\mathcal{I}_K(x),\\[4pt]
0,   & \text{otherwise},
\end{cases}
\]
where $\mathcal{I}_K(x)$ indexes the $K$ components of $x$ with largest absolute values (ties broken arbitrarily).  
The compressor satisfies the \emph{contraction property}
\[
\label{eq:contraction}
\E\!\bigl[\norm{\topk{x}-x}^2\bigr]\;\le\;(1-\delta)\,\norm{x}^2,\qquad 
\delta:=\frac{K}{d}\in(0,1].
\tag{C}
\]

\paragraph{Update rule.}
With learning rate $\eta>0$,
\[
c_t \;:=\; \topk{u_t},\qquad 
e_{t+1}\;:=\;u_t-c_t,\qquad 
w_{t+1}\;:=\;w_t-\eta\,c_t .
\tag{1}
\]

\section*{Pseudocode}
\begin{algorithm}[H]
\caption{Top‑$K$ Error‑Feedback SGD}
\begin{algorithmic}[1]
\State \textbf{Input:} initial point $w_0\in\R^d$, stepsize $\eta>0$, sparsity $K$.
\State Initialize error vector $e_0\gets 0$.
\For{$t=0,1,\dots,T-1$}
    \State Sample data $\xi_t$ and compute stochastic gradient $g_t\gets\nabla f(w_t;\xi_t)$.
    \State $u_t \gets g_t+e_t$ \Comment{error‑corrected gradient}
    \State $c_t \gets \topk{u_t}$ \Comment{keep $K$ largest entries}
    \State $e_{t+1} \gets u_t-c_t$ \Comment{update compression error}
    \State $w_{t+1} \gets w_t-\eta\,c_t$ \Comment{model update}
\EndFor
\State \textbf{Output:} $w_T$ (or averaged iterate $\bar w_T:=\frac1T\sum_{t=0}^{T-1}w_t$)
\end{algorithmic}
\end{algorithm}

\section*{Dry Run Example}
Consider the one‑dimensional quadratic $f(w)=(w-3)^2$, thus $\nabla f(w)=2(w-3)$.  
Take $w_0=0$, stepsize $\eta=0.1$, $K=1$ (no compression in 1‑D), and assume deterministic gradients (i.e.\ $\xi_t$ irrelevant).

\begin{center}
\begin{tabular}{c|c|c|c|c}
$t$ & $g_t$ & $e_t$ & $c_t$ & $w_{t+1}$ \\ \hline
0 & $2(0-3)=-6$ & $0$ & $-6$ & $0-0.1(-6)=0.6$ \\
1 & $2(0.6-3)=-4.8$ & $0$ & $-4.8$ & $0.6-0.1(-4.8)=1.08$ \\
2 & $2(1.08-3)=-3.84$ & $0$ & $-3.84$ & $1.08-0.1(-3.84)=1.464$ \\
\end{tabular}
\end{center}
The corresponding objective values are $f(w_1)=5.76$, $f(w_2)=4.6656$, $f(w_3)=3.7955$, showing descent toward the optimum $w^\star=3$.

\newpage
\section*{Assumptions}
\begin{assumption}[Smoothness]\label{as:smooth}
$f$ is $L$‑smooth: for all $x,y\in\R^d$,
\[
f(y)\le f(x)+\inner{\nabla f(x)}{y-x}+\frac{L}{2}\norm{y-x}^2 .
\]
\end{assumption}
\begin{assumption}[Unbiased stochastic gradients]\label{as:unbiased}
\[
\E\bigl[g_t\mid w_t\bigr]=\nabla f(w_t),\qquad 
\E\bigl[\norm{g_t-\nabla f(w_t)}^2\mid w_t\bigr]\le\sigma^2 .
\]
\end{assumption}
\begin{assumption}[Compression]\label{as:compress}
The top‑$K$ operator satisfies the contraction property \eqref{eq:contraction} with $\delta=K/d$.
\end{assumption}
\begin{assumption}[Bounded iterates]\label{as:bounded}
The level set $\{w\mid f(w)\le f(w_0)\}$ is bounded, guaranteeing $\norm{w_t}\le R$ for all $t$. (Used only to ensure existence of $f^\star:=\inf f$ and finiteness of constants.)
\end{assumption}

\section*{Main Convergence Theorem}
\begin{theorem}[Non‑convex convergence of EF‑TopK‑SGD]\label{thm:main}
Let Assumptions \ref{as:smooth}--\ref{as:compress} hold. Choose a constant stepsize
\[
0<\eta\le \frac{\delta}{L}.
\]
Then for any $T\ge1$ the iterates generated by \eqref{1} satisfy
\[
\frac{1}{T}\sum_{t=0}^{T-1}\E\!\bigl[\norm{\nabla f(w_t)}^2\bigr]
\;\le\;
\underbrace{\frac{2\bigl(f(w_0)-f^\star\bigr)}{\eta\delta\,T}}_{\text{optimization term}}
\;+\;
\underbrace{\frac{L\eta\sigma^2}{\delta}}_{\text{variance term}} .
\tag{2}
\]
Consequently, setting $\eta = \Theta\!\bigl(\frac{1}{\sqrt{T}}\bigr)$ yields
\[
\frac{1}{T}\sum_{t=0}^{T-1}\E\!\bigl[\norm{\nabla f(w_t)}^2\bigr]=
\mathcal{O}\!\Bigl(\frac{1}{\sqrt{T}}\Bigr) .
\]
\end{theorem}

\section*{Theoretical Properties}
\begin{enumerate}[label=\arabic*.]
\item \textbf{Stability w.r.t. compression.} The factor $\delta=K/d$ appears only as a multiplicative constant. As $K$ decreases, the bound deteriorates linearly, reflecting the loss of information.
\item \textbf{Hyper‑parameter sensitivity.} The stepsize must respect $\eta\le \delta/L$. Larger $K$ (larger $\delta$) permits a larger admissible $\eta$, matching intuition that more communicated information allows faster learning.
\item \textbf{Comparison to baselines.}
\begin{itemize}
\item \emph{Full‑gradient SGD} ($K=d$, $\delta=1$) reduces \eqref{2} to the classical non‑convex SGD bound.
\item \emph{SignSGD / Top‑1 EF} ($K=1$) incurs a factor $d$ penalty in the stepsize and variance term, reproducing known slower rates.
\end{itemize}
\item \textbf{Error‑feedback effect.} The error recursion guarantees that the compression bias does not accumulate; the error term $e_t$ remains bounded and vanishes in the average gradient bound.
\end{enumerate}

\section*{Discussion}
The theorem shows that despite aggressive sparsification, EF‑TopK‑SGD attains the same $\mathcal{O}(1/\sqrt{T})$ rate as vanilla SGD up to the $\delta^{-1}$ factor. This matches empirical observations that error feedback neutralises the bias introduced by top‑$K$ compression. The bound is tight in the sense that setting $K=d$ recovers the optimal constant. Moreover, the proof technique is modular: any unbiased compressor satisfying \eqref{eq:contraction} can replace top‑$K$, yielding identical rates with the appropriate $\delta$.

\newpage
\section*{Preliminaries (Definitions and Lemmas)}
\begin{lemma}[Descent Lemma for $L$‑smooth functions]\label{lem:descent}
For any $x,y\in\R^d$,
\[
f(y)\le f(x)+\inner{\nabla f(x)}{y-x}+\frac{L}{2}\norm{y-x}^2 .
\]
\end{lemma}
\begin{lemma}[Compression error recursion]\label{lem:error}
Define $e_{t+1}=u_t-\topk{u_t}$. Then for any $t$,
\[
\E\!\bigl[\norm{e_{t+1}}^2\mid \mathcal{F}_t\bigr]\le (1-\delta)\,\E\!\bigl[\norm{u_t}^2\mid\mathcal{F}_t\bigr] .
\tag{3}
\]
\end{lemma}
\begin{proof}
By the contraction property \eqref{eq:contraction} with $x=u_t$ and noting $e_{t+1}=u_t-\topk{u_t}$,
\[
\E\!\bigl[\norm{e_{t+1}}^2\mid\mathcal{F}_t\bigr]=
\E\!\bigl[\norm{u_t-\topk{u_t}}^2\mid\mathcal{F}_t\bigr]
\le (1-\delta)\norm{u_t}^2 .
\]
\end{proof}
Here $\mathcal{F}_t:=\sigma\{w_0,e_0,\dots,w_t,e_t,\xi_0,\dots,\xi_{t-1}\}$ denotes the natural filtration.

\section*{Main Proof (Step‑by‑Step Derivation)}
We prove Theorem \ref{thm:main}. Throughout the proof all expectations are taken w.r.t.\ the randomness up to iteration $t$ unless otherwise noted.

\bigskip
\noindent\textbf{[Step 1] Apply smoothness to the update.}  
Using Lemma \ref{lem:descent} with $x=w_t$, $y=w_{t+1}=w_t-\eta c_t$,
\begin{align}
f(w_{t+1}) &\le f(w_t)-\eta\inner{\nabla f(w_t)}{c_t}
               +\frac{L\eta^2}{2}\norm{c_t}^2 .
\tag{4}
\end{align}

\noindent\textbf{[Step 2] Decompose $c_t$ using the error‑feedback definition.}  
Recall $c_t=\topk{u_t}$ and $u_t=g_t+e_t$. Hence
\[
c_t = u_t - e_{t+1} .
\tag{5}
\]

\noindent\textbf{[Step 3] Expand the inner product in \eqref{4}.}  
Insert \eqref{5}:
\begin{align}
\inner{\nabla f(w_t)}{c_t}
&= \inner{\nabla f(w_t)}{u_t} - \inner{\nabla f(w_t)}{e_{t+1}} .
\tag{6}
\end{align}
Taking conditional expectation w.r.t.\ $\mathcal{F}_t$ and using $\E[g_t\mid\mathcal{F}_t]=\nabla f(w_t)$,
\[
\E\!\bigl[\inner{\nabla f(w_t)}{u_t}\mid\mathcal{F}_t\bigr]
= \norm{\nabla f(w_t)}^2 + \inner{\nabla f(w_t)}{e_t}.
\tag{7}
\]

\noindent\textbf{[Step 4] Bound the error term $\inner{\nabla f(w_t)}{e_t}$.}  
By Cauchy–Schwarz and Young’s inequality with parameter $\beta>0$,
\[
\inner{\nabla f(w_t)}{e_t}\le \frac{1}{2\beta}\norm{\nabla f(w_t)}^2
                         +\frac{\beta}{2}\norm{e_t}^2 .
\tag{8}
\]

\noindent\textbf{[Step 5] Take expectation of the descent inequality.}  
From \eqref{4}–\eqref{8},
\begin{align}
\E\!\bigl[f(w_{t+1})\mid\mathcal{F}_t\bigr]
&\le f(w_t)-\eta\Bigl(\norm{\nabla f(w_t)}^2
      +\inner{\nabla f(w_t)}{e_t}
      -\inner{\nabla f(w_t)}{e_{t+1}}\Bigr)
   +\frac{L\eta^2}{2}\E\!\bigl[\norm{c_t}^2\mid\mathcal{F}_t\bigr] .
\tag{9}
\end{align}
Insert the bound \eqref{8} for $\inner{\nabla f(w_t)}{e_t}$ and keep $\inner{\nabla f(w_t)}{e_{t+1}}$ as is:
\begin{align}
\E\!\bigl[f(w_{t+1})\mid\mathcal{F}_t\bigr]
&\le f(w_t)-\eta\Bigl(1-\tfrac{1}{2\beta}\Bigr)\norm{\nabla f(w_t)}^2
      -\eta\inner{\nabla f(w_t)}{e_{t+1}}
      +\frac{\eta\beta}{2}\norm{e_t}^2
      +\frac{L\eta^2}{2}\E\!\bigl[\norm{c_t}^2\mid\mathcal{F}_t\bigr] .
\tag{10}
\end{align}

\noindent\textbf{[Step 6] Relate $\norm{c_t}^2$ to $\norm{u_t}^2$.}  
From \eqref{5},
\[
\norm{c_t}^2 = \norm{u_t-e_{t+1}}^2
            = \norm{u_t}^2 + \norm{e_{t+1}}^2 - 2\inner{u_t}{e_{t+1}} .
\]
Taking conditional expectation and using $\E[e_{t+1}\mid\mathcal{F}_t]=0$ (since compression is unbiased in expectation),
\[
\E\!\bigl[\norm{c_t}^2\mid\mathcal{F}_t\bigr]
= \E\!\bigl[\norm{u_t}^2\mid\mathcal{F}_t\bigr]
  +\E\!\bigl[\norm{e_{t+1}}^2\mid\mathcal{F}_t\bigr] .
\tag{11}
\]

\noindent\textbf{[Step 7] Upper‑bound $\E[\norm{u_t}^2]$.}  
Recall $u_t=g_t+e_t$. Using unbiasedness and variance bound,
\begin{align}
\E\!\bigl[\norm{u_t}^2\mid\mathcal{F}_t\bigr]
&= \E\!\bigl[\norm{g_t+e_t}^2\mid\mathcal{F}_t\bigr] \nonumber\\
&= \norm{e_t}^2 + \norm{\nabla f(w_t)}^2
   + \E\!\bigl[\norm{g_t-\nabla f(w_t)}^2\mid\mathcal{F}_t\bigr]
   + 2\inner{\nabla f(w_t)}{e_t} \nonumber\\
&\le \norm{e_t}^2 + \norm{\nabla f(w_t)}^2 + \sigma^2
   + \frac{1}{\beta}\norm{\nabla f(w_t)}^2
   + \beta \norm{e_t}^2 \qquad\text{(by Young's inequality)} \nonumber\\
&= \Bigl(1+\frac{1}{\beta}\Bigr)\norm{\nabla f(w_t)}^2
   + (1+\beta)\norm{e_t}^2 + \sigma^2 .
\tag{12}
\end{align}

\noindent\textbf{[Step 8] Bound $\E[\norm{e_{t+1}}^2]$ using Lemma \ref{lem:error}.}  
From \eqref{3},
\[
\E\!\bigl[\norm{e_{t+1}}^2\mid\mathcal{F}_t\bigr]\le (1-\delta)\,\E\!\bigl[\norm{u_t}^2\mid\mathcal{F}_t\bigr] .
\tag{13}
\]

\noindent\textbf{[Step 9] Combine (11)–(13).}  
Plug \eqref{13} into \eqref{11} and then use \eqref{12}:
\begin{align}
\E\!\bigl[\norm{c_t}^2\mid\mathcal{F}_t\bigr]
&\le \E\!\bigl[\norm{u_t}^2\mid\mathcal{F}_t\bigr]
   + (1-\delta)\E\!\bigl[\norm{u_t}^2\mid\mathcal{F}_t\bigr] \nonumber\\
&= (2-\delta)\E\!\bigl[\norm{u_t}^2\mid\mathcal{F}_t\bigr] \nonumber\\
&\le (2-\delta)\Bigl[\Bigl(1+\frac{1}{\beta}\Bigr)\norm{\nabla f(w_t)}^2
   + (1+\beta)\norm{e_t}^2 + \sigma^2\Bigr] .
\tag{14}
\end{align}

\noindent\textbf{[Step 10] Substitute (14) into (10).}  
We obtain
\begin{align}
\E\!\bigl[f(w_{t+1})\mid\mathcal{F}_t\bigr]
&\le f(w_t)-\eta\Bigl(1-\tfrac{1}{2\beta}\Bigr)\norm{\nabla f(w_t)}^2
   -\eta\inner{\nabla f(w_t)}{e_{t+1}}
   +\frac{\eta\beta}{2}\norm{e_t}^2  \nonumber\\
&\quad +\frac{L\eta^2}{2}(2-\delta)
   \Bigl[\Bigl(1+\frac{1}{\beta}\Bigr)\norm{\nabla f(w_t)}^2
   + (1+\beta)\norm{e_t}^2 + \sigma^2\Bigr] .
\tag{15}
\end{align}

\noindent\textbf{[Step 11] Choose $\beta$ to cancel the mixed term $\inner{\nabla f(w_t)}{e_{t+1}}$.}  
Taking full expectation and noting $\E[\inner{\nabla f(w_t)}{e_{t+1}}]=0$ (since $e_{t+1}$ is zero‑mean conditioned on $\mathcal{F}_t$), the term disappears.  

Now set $\beta = \frac{1}{2}$ (any constant $>0$ works; this choice simplifies constants). Then
\[
1-\frac{1}{2\beta}=0,\qquad
1+\frac{1}{\beta}=3,\qquad
1+\beta=\frac{3}{2}.
\]

\noindent\textbf{[Step 12] Simplify (15) with $\beta=\tfrac12$.}  
\begin{align}
\E\!\bigl[f(w_{t+1})\bigr]
&\le \E\!\bigl[f(w_t)\bigr]
   +\frac{L\eta^2}{2}(2-\delta)
   \Bigl[3\,\norm{\nabla f(w_t)}^2
   +\tfrac{3}{2}\norm{e_t}^2 + \sigma^2\Bigr]
   +\frac{\eta}{4}\norm{e_t}^2 .
\tag{16}
\end{align}
Re‑arrange to isolate the gradient term:
\[
\E\!\bigl[f(w_t)-f(w_{t+1})\bigr]
\ge \underbrace{\Bigl( \eta - \tfrac{3L\eta^2}{2}(2-\delta) \Bigr)}_{A}
        \norm{\nabla f(w_t)}^2
     -\underbrace{\Bigl(\tfrac{\eta}{4}
        +\tfrac{3L\eta^2}{4}(2-\delta)\Bigr)}_{B}\norm{e_t}^2
     -\tfrac{L\eta^2}{2}(2-\delta)\sigma^2 .
\tag{17}
\]

\noindent\textbf{[Step 13] Ensure $A>0$ by step‑size restriction.}  
Require
\[
\eta \le \frac{2}{3L(2-\delta)} .
\tag{18}
\]
Since $2-\delta\le 2$, a sufficient condition is $\eta\le\frac{\delta}{L}$ (used in the theorem statement). Under this condition $A\ge \frac{\eta\delta}{2}$.

\noindent\textbf{[Step 14] Control the error term $ \norm{e_t}^2$ .}  
From Lemma \ref{lem:error} and (12) with $\beta=\frac12$,
\begin{align}
\E\!\bigl[\norm{e_{t+1}}^2\bigr]
&\le (1-\delta)\E\!\bigl[\norm{u_t}^2\bigr] \nonumber\\
&\le (1-\delta)\Bigl[3\norm{\nabla f(w_t)}^2
      +\tfrac{3}{2}\norm{e_t}^2 + \sigma^2\Bigr] .
\tag{19}
\end{align}
Iterating (19) yields a geometric series; in particular $ \norm{e_t}^2$ remains bounded by a constant multiple of $\sigma^2/\delta$. For the purpose of the average bound we can drop the non‑negative term $B\norm{e_t}^2$ (it only makes the RHS smaller).

\noindent\textbf{[Step 15] Summation over $t=0,\dots,T-1$.}  
Sum \eqref{17} and discard the $B\norm{e_t}^2$ term:
\[
\sum_{t=0}^{T-1}\E\!\bigl[f(w_t)-f(w_{t+1})\bigr]
\ge \frac{\eta\delta}{2}\sum_{t=0}^{T-1}\E\!\bigl[\norm{\nabla f(w_t)}^2\big]
   -\frac{L\eta^2}{2}(2-\delta)\sigma^2 T .
\]
The left‑hand side telescopes to $f(w_0)-\E[f(w_T)]\le f(w_0)-f^\star$. Hence
\[
\frac{\eta\delta}{2}\sum_{t=0}^{T-1}\E\!\bigl[\norm{\nabla f(w_t)}^2\big]
\le f(w_0)-f^\star +\frac{L\eta^2}{2}(2-\delta)\sigma^2 T .
\tag{20}
\]

\noindent\textbf{[Step 16] Divide by $T$ and rearrange.}  
\[
\frac{1}{T}\sum_{t=0}^{T-1}\E\!\bigl[\norm{\nabla f(w_t)}^2\big]
\le
\frac{2\bigl(f(w_0)-f^\star\bigr)}{\eta\delta\,T}
   +\frac{L\eta(2-\delta)}{\delta}\,\sigma^2 .
\tag{21}
\]
Since $2-\delta\le 2$, we further bound the variance term by $\frac{L\eta\sigma^2}{\delta}$, obtaining exactly the bound \eqref{2} claimed in Theorem \ref{thm:main}.

\noindent\textbf{[Step 17] Choice of stepsize for optimal rate.}  
Set $\eta = \frac{c}{\sqrt{T}}$ with $c\le \frac{\delta}{L}$; then the first term in \eqref{2} scales as $O(1/\sqrt{T})$ and the second term as $O(1/\sqrt{T})$, establishing the $\mathcal{O}(1/\sqrt{T})$ rate.

\qed

\section*{Rate Derivation (With Constants)}
From \eqref{21} with $\eta = \frac{\delta}{L\sqrt{T}}$ (which respects $\eta\le\delta/L$):
\[
\frac{1}{T}\sum_{t=0}^{T-1}\E\!\bigl[\norm{\nabla f(w_t)}^2\bigr]
\le
\underbrace{\frac{2L\bigl(f(w_0)-f^\star\bigr)}{\delta^2}}_{\displaystyle C_1}\frac{1}{\sqrt{T}}
\;+\;
\underbrace{\frac{\delta}{L}\sigma^2}_{\displaystyle C_2}\frac{1}{\sqrt{T}} .
\]
Thus the explicit constant is $C = C_1+C_2$.

\section*{Special Cases}
\begin{enumerate}[label=\alph*)]
\item \textbf{No compression ($K=d$).} Then $\delta=1$ and the bound reduces to
\[
\frac{1}{T}\sum_{t=0}^{T-1}\E\!\bigl[\norm{\nabla f(w_t)}^2\bigr]
\le
\frac{2\bigl(f(w_0)-f^\star\bigr)}{\eta T}
   + L\eta\sigma^2 ,
\]
which is the classical SGD guarantee.
\item \textbf{Extreme sparsification ($K=1$).} Here $\delta=1/d$; the stepsize must be at most $1/(dL)$. The bound deteriorates by a factor $d$, matching known lower bounds for 1‑bit/SignSGD‑type methods.
\item \textbf{Deterministic gradients ($\sigma=0$).} The variance term vanishes, yielding a pure optimization term $O(1/T)$, i.e.\ linear convergence of the gradient norm to zero for smooth non‑convex objectives under a diminishing stepsize.
\end{enumerate}

\end{document}